\documentclass[10pt,a4paper]{amsart}
\usepackage[margin=19mm]{geometry}
\usepackage{amsmath,amssymb,mathtools}
\usepackage[scr=rsfs,cal=euler]{mathalfa}
\usepackage{xfrac}
\usepackage[unicode,hidelinks,bookmarks=false]{hyperref}
\newcommand{\ie}{i.e.\ }
\newcommand{\cf}{cf.\ }
\newcommand{\resp}{respectively\ }
\newcommand{\Subsection}{\subsection}
\providecommand{\nobreakdash}{\nobreak\hbox{-}}
\setlength{\emergencystretch}{2em}
\multlinegap0pt
\usepackage{amssymb}
\usepackage{mathtools}
\usepackage{tensor}
\newtheorem{definition}{Definition}
\newtheorem{lemma}{Lemma}
\newcommand{\N}{\mathbb{N}}
\title{Compact Quantum Metric Spaces from Weakly Geodesic Length Functions on Hyperbolic Groups}
\author{Are Austad}
\date{Source: arXiv:2608.28303v1 (CC BY 4.0)}
\begin{document}
\maketitle

\noindent\emph{Source and licence.} Compact Quantum Metric Spaces from Weakly Geodesic Length Functions on Hyperbolic Groups. Version arXiv:2608.28303v1 (CC BY 4.0), distributed by arXiv under the Creative Commons Attribution 4.0 International licence, http://creativecommons.org/licenses/by/4.0/, as recorded in the arXiv licence metadata for this exact version and shown on the abstract page https://arxiv.org/abs/2608.28303v1. Full licence notice: \texttt{licenses/arxiv-2608.28303-CC-BY-4.0.txt}.\par\smallskip

\noindent\textit{Editorial scope.} Counted unit: the article's lemma lemma:annular-norm-control (source0.tex lines 366--378). The article's complete original proof of it is delivered with it (source0.tex lines 379--558, one \texttt{proof} environment of 180 lines). No mathematics is written, repaired or re-derived: the statement and the proof are the article's own lines, in the article's own notation, and no retained line is substituted or re-flowed.  Retained uncounted as the source-local context the proof needs: lines 211--216.  Nothing was omitted from any retained span, so the package carries no omission record.  No retained line contains a citation, so the wrapper carries no bibliography and no bibliography entry.  No retained line contains a figure environment, an image-inclusion command or drawing code, so no image is delivered and the package declares no asset dependency.  Editorial formatting only, in addition: an AMS \texttt{amsart} preamble that restates the article's own declarations and macros used by the retained lines, namely the environments \texttt{definition}, \texttt{lemma} on separate counters, together with the standard packages those lines need.  The retained environments print in this document as Definition 1, Lemma 1 in order of appearance, whereas the article prints the same items with the numbering of its own full document.  The complete native source of the article as distributed in the arXiv e-print for this exact version is bundled unchanged as \texttt{sources/208779/source0.tex}.
\begin{definition}\label{def:hyperbolic-metric-space}
    A metric space $(X,d)$ is said to be \emph{hyperbolic} (or $\delta$-hyperbolic) if there is $\delta \geq 0$ such that for any four points $x,y,z,w \in X$ we have
    \begin{align*}
        d(x,y) + d(z,w) \leq \max \{ d(x,z) + d(y,w), d(x,w) + d(y,z) \} + \delta .
    \end{align*}
\end{definition}

\begin{lemma}\label{lemma:annular-norm-control}
   Let $G$ be a 
   countable discrete
   group, and suppose $\ell \colon G \to [0,\infty)$ is a proper length function such that $(G,d_\ell)$ is a $\delta$-hyperbolic and weakly $c$-geodesic metric space. 
   Denote by $E_m := \{x \in G \mid m \leq \ell(x) < m+1 \}$, 
   and let 
   $P_m \in B(\ell^2(G))$ be the orthogonal projection onto $\ell^2(E_m)$. 
   Then there exists $C \geq 0$, depending only on $G$ and $\ell$, such that for all $m,n,k \in \N_0$ and any $f \in C_c(G)$ supported on $E_k$, we have
   \begin{align*}
       \Vert P_m \lambda(f) P_n \Vert \leq C \Vert f \Vert_2. 
   \end{align*}
   Explicitly, we may take any $C \geq |B_\ell(2c + \frac{7}{2}+2\delta)\vert$. 
\end{lemma}

\begin{proof}
    We need to show that
    \begin{align*}
        \Vert P_m(f * \xi) \Vert \leq C \Vert f\Vert_2 \Vert \xi \Vert_2
    \end{align*}
    for any $\xi \in \ell^2(G)$  supported on $E_n$. Fix such a $\xi$, and set $X_{m,n,k} := E_m \cap E_kE_n$. If $x \in E_m \setminus X_{m,n,k}$, then $f*\xi(x) = 0$, and so it suffices to consider $x \in X_{m,n,k}$. For such $x$, let $y \in E_k$ and $z \in E_n$ be arbitrary elements such that $yz = x$. 

    Now set
    \begin{align*}
        a_x = \frac{k+n+1 - \ell(x)}{2} , \quad b_x = \frac{k-n +\ell(x)}{2} .
    \end{align*}
     Since $\ell(x) \geq \ell(z) - \ell(y) > n - (k+1)$ for $y \in E_k$ and $z \in E_n$, we see that $b_x > -\frac{1}{2}$. 
    Using $\ell(x) \geq \ell(y) - \ell(z)$ we also establish that $\ell(x) - b_x > -\frac{1}{2}$. 
    
    Now define
    \begin{align*}
        r_x = \begin{cases}
            0 & b_x < 0 \\
            b_x & 0\leq b_x \leq \ell(x) \\
            \ell(x) & b_x > \ell(x)
        \end{cases}
    \end{align*}
    which immediately implies $|r_x - b_x| < \frac{1}{2}$ by the above observations. 
    For each $x \in X_{m,n,k}$, we use weak geodesicity to fix an element $\overline{x} \in G$ satisfying 
    \begin{align*}
        |\ell(\overline{x}) - r_x | \leq c, \quad | \ell(\overline{x}^{-1}x) - (\ell(x) - r_x)| \leq c,
    \end{align*}
    and set $\tilde{x} := \overline{x}^{-1}x$.

    Applying the hyperbolicity criterion from Definition \ref{def:hyperbolic-metric-space} to the points $e, x, \overline{x}, y$, we obtain
    \begin{align*}
        \ell(x) + \ell(y^{-1}\overline{x}) \leq \max \{ \ell(\overline{x}) + \ell(z) , \ell(y) + \ell(\tilde{x}) \} + \delta, 
    \end{align*}
    where we have used $z = y^{-1}x$ and $\tilde{x} = \overline{x}^{-1}x$. Now,
    \begin{align*}
        \ell(\overline{x}) + \ell(z) < r_x + c  + n +1 \leq b_x + \frac{1}{2} + c + n+1 = a_x + \ell(x) + c + 1.
    \end{align*}
    Moreover,
    \begin{align*}
        \ell(y) + \ell(\overline{x}^{-1}x) < k+1 + \ell(x) - r_x +c \leq k+1 +\ell(x) -b_x + \frac{1}{2} +c  = a_x + \ell(x) +c + 1.
    \end{align*}
    This now implies
    \begin{align*}
        \ell(y^{-1}\overline{x}) \leq - \ell(x) + a_x + \ell(x) + c + 1 + \delta = a_x + c +1+\delta
    \end{align*}
    Set $R_x := a_x + c+ 1 +\delta$. 
    Since the preceding estimate holds for every factorization $x = yz$ with $y \in E_k$ and $z \in E_n$, the Cauchy--Schwarz inequality yields
    \begin{align*}
        \Vert P_m(f*\xi)\Vert_2^2
\leq
\sum_{x\in X_{m,n,k}}
\sum_{\ell(u)\leq R_x}
\sum_{\ell(v)\leq R_x}
|f(\bar xu)|^2|\xi(v\widetilde x)|^2.
    \end{align*}

    For each term in the triple sum, define $y = \bar{x}u$ and $z = v\tilde{x}$. We only have nonzero contributions for $y \in E_k$ and $z \in E_n$. 
    We can then write
    \begin{align}\label{eq:multiplicity-convolution-bound-Nyz}
        \Vert P_m (f*\xi) \Vert^2 \leq  \sum_{y \in E_k} \sum_{z \in E_n} N(y,z) \vert f(y) \vert^2 \cdot \vert \xi(z) \vert^2
    \end{align}
    for
    \begin{align*}
        N(y,z) := |\{(x,u,v) \mid x \in X_{m,n,k}, y = \bar{x}u, z= v\tilde{x}\text{ and } \ell(u),\ell(v) \leq R_x \}|.
    \end{align*}
    It suffices to show that $N(y,z)$ is uniformly bounded for all $y$ and $z$. To this end, we define
    \begin{align*}
        \overline{F}_y :&= \{ u \in G \mid  \text{  there is $x \in X_{m,n,k}$ such that $y = \overline{x}u$} \text{ and } \ell(u) \leq R_x \} \\
        \tilde{F}_z :&=  \{ v \in G \mid  \text{  there is $x \in X_{m,n,k}$ such that $z = v\tilde{x}$ and $\ell(v) \leq R_x$} \}.
    \end{align*}
    For fixed $y,z,u,v$, the element $x$ is uniquely determined by $x=\bar x\widetilde x
=yu^{-1}v^{-1}z$.
Hence the map
$(x,u,v)\mapsto (u,v)$ 
is injective on the set counted by $N(y,z)$, and therefore
\begin{align*}
    N(y,z)\leq |F_y|\,|\widetilde F_z|.
\end{align*}  
    We therefore aim to show that $\vert \overline{F}_y \vert $ and $\vert \tilde{F}_z \vert$ are bounded independently of $y$ and $z$. 


    Let $u,w \in \overline{F}_y$, and choose $x,x' \in X_{m,n,k}$ such that
    \begin{align*}
        y = \overline{x}u = \overline{x'}w
    \end{align*}
    with $\ell(u) \leq R_x$ and $\ell(w) \leq R_{x'}$. Set $s:= \overline{x}$ and $t:= \overline{x'}$, so that $y = su = tw$. Since $x,x' \in E_m$, we have
    $\vert \ell(x) - \ell(x') \vert < 1$ and therefore $\vert b_x - b_{x'} \vert \leq \frac{1}{2}$. Moreover, because $x,x' \in X_{m,n,k}$, we have
    \begin{align*}
        \vert r_x - b_x \vert < \frac{1}{2} \quad \text{and} \quad \vert r_{x'} - b_{x'} \vert < \frac{1}{2}.
    \end{align*}
    It follows that $\vert r_x - r_{x'}\vert < \frac{3}{2}$ by the triangle inequality. 
    
    By construction we have
    \begin{align*}
        \ell(s) &\leq r_{x} +c, \quad \ell(u) \leq a_x + 1 +c +\delta \\
        \ell(t) &\leq r_{x'} + c, \quad \ell(w) \leq a_{x'} +1+c+\delta.
    \end{align*}
    Now, $a_x = k - b_x + \frac{1}{2}$ and $|b_x - r_x | < \frac{1}{2}$, and similarly for $x'$. Thus
    \begin{align*}
        a_x \leq k+1 - r_x , \quad a_{x'} \leq k+1 - r_{x'},
    \end{align*}
    and so 
    \begin{align*}
        \ell(u) &\leq k-r_x +2+c+\delta \\ 
        \ell(w) &\leq k-r_{x'} +2 +c +\delta.
    \end{align*}
    In particular,
    \begin{align*}
        \ell(s) + \ell(w) &\leq r_x + c + k - r_{x'} +2 +c +\delta \leq k + 2c + \frac{7}{2} + \delta \\
        \ell(t) + \ell(u) &\leq r_{x'} + c + k - r_{x} +2+c+\delta \leq k + 2c  + \frac{7}{2} + \delta
    \end{align*}
    where we have used $|r_x - r_{x'}| < \frac{3}{2}$. 
    We then apply the hyperbolicity condition to $e,y,s,t$ to obtain
    \begin{align*}
        \ell(y) + \ell(s^{-1}t) \leq \max \{ \ell(s) + \ell(w), \ell(t) + \ell(u) \} +\delta ,
    \end{align*}
    and since $\ell(y) \geq k$, we deduce
    \begin{align*}
        \ell(u w^{-1}) = \ell(s^{-1}t) \leq 2c  + \frac{7}{2} + 2\delta.
    \end{align*}

    If $\overline{F}_y = \emptyset$, the bound is immediate. Otherwise,
    fix $w \in \overline{F}_y$. The map $\overline{F}_y \to G$ , $u \mapsto uw^{-1}$
    is injective and takes values in $B_\ell(2c +\frac{7}{2} + 2\delta )$. Thus
    \begin{align*}
        \vert \overline{F}_y \vert \leq | B_\ell(2c +\frac{7}{2} + 2\delta ) \vert.
    \end{align*}
    Since $\ell$ is proper, this is a finite value, independent of $y$. 

    The same bound holds for $\tilde{F}_z$. Indeed, 
$\tilde{x}=\bar x^{-1}x$, where $\bar x$ was chosen so that
\begin{align*}
    |\ell(\bar x)-r_x|\leq c,
    \qquad
    |\ell(\tilde{x})-(\ell(x)-r_x)|\leq c.
\end{align*}
It follows that $\tilde{x}^{-1}$ is a $c$-approximate intermediate
point for $x^{-1}$ at distance $\ell(x)-r_x$ from $e$. Indeed, since $\ell(\tilde{x}^{-1}) = \ell(\tilde{x})$ we have
\begin{align*}
    |\ell(\tilde{x}^{-1})-(\ell(x)-r_x)|\leq c,
\end{align*}
while
$(\tilde{x}^{-1})^{-1}x^{-1}
    =\tilde{x} x^{-1}
    =\bar x^{-1}$,
and thus
\begin{align*}
     \bigl|\ell((\tilde{x}^{-1})^{-1}x^{-1})-r_x\bigr|
    =|\ell(\bar x)-r_x|
    \leq c.
\end{align*} 
These are precisely the two inequalities
for an approximate intermediate point for $x^{-1}$ corresponding to
the 
parameter $\ell(x)-r_x$.

Now, if $v\in\tilde{F}_z$, then for some $x\in X_{m,n,k}$, we have $z=v\tilde{x}$ with $\ell(v)\leq R_x$, 
and therefore
\begin{align*}
    z^{-1}=\tilde{x}^{-1}v^{-1}.
\end{align*}
Noting that under $x \mapsto x^{-1}$ and the interchange of $k$ and $n$, $a_x$ remains unchanged and the corresponding $r$-parameter is $\ell(x) - r_x$, the preceding argument for $\overline{F}_y$ applied mutatis mutandis to $z^{-1}$  bounds the number of possible
values of $v^{-1}$ by $|B_\ell\!\left(2c+\frac72+2\delta\right)|$. 
Since inversion is a bijection, we conclude that
\begin{align*}
    |\tilde{F}_z|
    \leq
    \left|B_\ell\!\left(2c+\frac72+2\delta\right)\right|.
\end{align*}
    Set 
    \begin{align*}
        C := \vert B_\ell(2c + \frac{7}{2} + 2\delta) \vert.
    \end{align*}
    
    By \eqref{eq:multiplicity-convolution-bound-Nyz} above we immediately see    
    \begin{align*}
         \Vert P_m (f*\xi) \Vert^2 \leq  \sum_{y \in E_k} \sum_{z\in E_n} N(y,z) \vert f(y) \vert^2 \cdot \vert \xi(z) \vert^2 \leq C^2 \bigg( \sum_{y\in E_k} \vert f(y) \vert^2 \bigg)\bigg( \sum_{z \in E_n} \vert \xi(z) \vert^2 \bigg),
    \end{align*}
    from which we deduce $\Vert P_m (f*\xi)\Vert \leq C \Vert f\Vert_2 \Vert \xi \Vert_2$ and thus $\Vert P_m \lambda(f) P_n \Vert \leq C \Vert f \Vert_2$. This finishes the proof. 
\end{proof}

\end{document}
